\section{Introduction}
\label{sec:intro}

Parameterised quantum circuits trained by gradient descent are the most common model in near-term
quantum machine learning \citep{preskill2018nisq,cerezo2021vqa}. Software such as TensorFlow Quantum
(TFQ) \citep{broughton2020tfq} lets anyone train them on a laptop. That makes it easy to run a demo, and
also easy to misread one: without a matched classical baseline, a sense of what the data allows, or a
look at gradients and noise, a working QNN says little.

This paper is a compact, reproducible tour of those checks on one task: the Farhi--Neven QNN
\citep{farhi2018classification} on $4\times4$ MNIST 3-vs-6, the setting of the TFQ tutorial. We contribute:
\begin{itemize}
  \item a matched comparison of the QNN, a classical network of similar size, and a hybrid model that
        adds a classical head to the same circuit, alongside the accuracy ceiling the input allows
        (\S\ref{sec:res-mnist}, \S\ref{sec:res-hybrid});
  \item measurements of gradient variance in random circuits as width and depth grow
        (\S\ref{sec:res-barren});
  \item the effect of depolarizing noise, both on trained models and during training
        (\S\ref{sec:res-noise});
  \item code that regenerates every figure, table and number from raw results with one command,
        running on current TFQ (\valVerTFQ) in Docker.
\end{itemize}
None of these phenomena are new. The point is to reproduce them together, with error bars, in a form
that is easy to rerun and extend.
